% ============================================================
% The Decision to Fight
% ============================================================

\section{The Decision to Fight}
\label{sec:decision}

Measure in hand, I now study whether intervention expectations affect the
decision to initiate civil war.  The \textit{Entrants} model augments a standard
onset logit with the two shadow variables constructed in the previous
section:

\begin{equation*}
U_O(F; \hat{\sigma}) = \mathbf{X}\boldsymbol{\beta} + \gamma^O \sum_i \hat{\sigma}_i(O) + \gamma^G \sum_i \hat{\sigma}_i(G) + \epsilon,
\end{equation*}

\noindent where $\mathbf{X}$ contains the baseline country-year predictors and the two
summations are the expected total number of opposition-biased and
government-biased interveners, respectively.\footnote{In estimation, the
summations enter through the inverse hyperbolic sine transform,
$\operatorname{asinh}(\sum_i \hat{\sigma}_i)$, for variance stabilization.  This
introduces a concavity---diminishing marginal returns to additional
expected interveners---that is not implied by the theory.  Because
$\operatorname{asinh}$ is strictly increasing, the sign of the estimated coefficients
is invariant to this choice, and the out-of-sample model comparisons
test whether the shadow variables improve predictions regardless of
functional form.  The exact shape of the utility function over
aggregated intervention expectations is not identified from these data.}
This is the simplest
specification of Equation~\eqref{eq:EU}, in which all potential interveners carry the same
utility weight ($u_O^i(O) = \gamma^O$ and $u_O^i(G) = \gamma^G$ for
all $i$).  I relax this assumption in the heterogeneous-utilities
extensions below.

The two-stage design extends the logic of statistical backwards induction
\citep{bas2008} to an $n$-player environment.  In the two-player case, backwards
induction estimates the second mover's choice probabilities and plugs them
into the first mover's expected utility.  Here the intervention subgame
involves up to 193 simultaneous movers; separability and the Nash
fixed-point reduce the problem to one in which each intervener's marginal
probabilities suffice, and Stage~2 recovers the opposition's utility
parameters from their aggregate effect on onset.

\subsection{Empirical Strategy}

Testing these models requires a comparison point.  Following
\citet{clarke2007}, the relevant question is not whether intervention
expectations matter in isolation but whether a model that includes them
outperforms one that does not.

The \textit{Baseline} model serves as the rival.  It is a logistic regression
for civil war onset using the covariate set from
\citet{fearon2003} Table~1, Model~1---the structural
specification the literature has converged on as a reference: prior war,
log per capita income, log population, mountainous terrain, oil exporter
status, new state status, political instability, ethnic and religious
fractionalization, and democracy---estimated on our sample (COW Intra-State Wars v5.1, 1946--2014).
This is not a replication of F\&L; it is their specification used as a
comparison point on our data with updated sources.  Time-varying
covariates (income, population, democracy, political instability) are
lagged one year to avoid post-onset contamination; income and population
come from V-Dem v15 \citep{vdem2025} and democracy from the Polity project.
Civil war history is drawn from COW v5.1.  Slowly changing or fixed
characteristics---terrain, ethnic and religious fractionalization,
colonial history---are taken from the F\&L replication dataset.  To account for the mechanical growth of the summation-based
shadow variables as the state system expands, I add a year trend to all
specifications.

\citet{cederman2010} challenge the aggregate-fractionalization
baseline on the grounds that ethnopolitical exclusion from state power is
the theoretically appropriate predictor.  The Coethnics extended
specification (Section~\ref{sec:decision}) partially addresses this critique by
interacting intervention expectations with ethnic matching between the
intervener and the host country's population.

Model comparison draws on both in-sample fit and out-of-sample
predictive performance.  Within sample, I report the proportional
reduction in log-loss relative to random guessing and the area under the
ROC curve.  Out-of-sample performance uses leave-one-war-out
cross-validation: all country-years belonging to a civil war are held
out together, and predictions are formed from the remaining data.  This
scheme respects temporal dependence and avoids the information leakage
that would arise from splitting individual country-years at random.

Because the Stage~1 probabilities are estimated quantities, I retain all
25 multiply-imputed shadow measures (5 country-year imputations $\times$ 5
directed-dyad imputations) for the uncertainty analysis described below.
Point estimates for model-fit metrics are computed on the averaged shadow;
coefficient standard errors account for the full spread of measurement
draws.

\subsection{Intervention Expectations and Onset}

Table~\ref{tab:coefs} presents the paper's central result.  The two
shadow variables enter with the signs the theory requires: expected
government-biased intervention is negative ($\gamma^G = -1.62$),
consistent with deterrence, and expected opposition-biased intervention
is positive ($\gamma^O = +0.64$), consistent with emboldening.  This
directional separation holds in every one of the 25 measurement draws,
in all ten specifications reported in Table~\ref{tab:topfit}, and in the
country fixed-effects estimator.

The Entrants specification improves on the Baseline both in-sample (PRL
from 10.1\% to 11.9\%, AUC from 0.771 to 0.787) and out-of-sample
(OOS PRL = $+0.4\%$, OOS AUC from 0.650 to 0.672).  For a rare-event
model with a 2\% base rate, positive OOS PRL is a demanding bar: the
log-loss metric penalizes confident false positives severely, and the
Baseline specification---which encodes decades of accumulated knowledge
about structural onset risk---does not clear it.

\subsubsection*{Generated-Regressor Uncertainty}

Because the shadow variables are predicted from a Stage~1 classifier
rather than directly observed, standard logistic-regression errors
understate uncertainty.  \citet{knoxlucascho2022} show that this is the
most common methodological failure in learned-proxy applications: of the
48 studies they review, nearly all ignore measurement-stage variability.
I follow their recommended $T \times P$ correction: the primary analysis
runs separately for each of the $T = 25$ shadow measures, each is
bootstrapped $P = 200$ times (pairs-cluster, resampling countries), and
the resulting 5,000 coefficient vectors are pooled.

The corrected standard errors in Table~\ref{tab:coefs} are roughly
$3.5\times$ larger for $\hat{E}^G$ and $2.9\times$ larger for
$\hat{E}^O$ than naive MLE output.  A variance decomposition shows that
88--92\% of this total variance originates in the measurement
stage---variation \textit{between} shadow draws---with only 8--12\% from
primary-analysis sampling.  The measurement model, not the regression, is
the binding source of uncertainty.

Two features of the data make this unsurprising.  First, 184 onsets
across 25 measurement draws leave the Stage~2 signal thin relative to
variation in the learned proxy.  Second, $\hat{E}^G$ and $\hat{E}^O$ are
correlated at $r = 0.94$: countries that attract government-biased
intervention expectations also attract opposition-biased expectations.
This collinearity means that small perturbations in the Stage~1
classifier move both shadow variables together, inflating the
between-draw variance of each coefficient.  Tested jointly, however, the
shadow variables are unambiguous: a likelihood-ratio test rejects the
Baseline in favor of Entrants at $p < 0.001$ (LR $= 32.7$, df $= 2$),
and most type-disaggregated specifications reject Entrants in turn
(Appendix~\ref{app:lr-tests}, Table~\ref{tab:lr-tests}).  The
collinearity inflates \textit{individual} coefficient uncertainty but
does not prevent the pair from carrying strong joint signal.  The
consistent sign separation across all 25 draws---$\gamma^G < 0$ and
$\gamma^O > 0$ in every one---reinforces this: the logit reliably
separates the two effects despite having only $\approx$6\% independent
variation to work with.

Neither coefficient achieves conventional significance individually: the
95\% confidence interval for $\gamma^G$ runs from $-4.63$ to $+0.69$
and for $\gamma^O$ from $-1.94$ to $+2.59$.  Under the average
monotonicity condition of \citet{knoxlucascho2022} (\S3.1)---more of the
true intervention expectation implies more of the proxy, which the
Stage~1 calibration diagrams verify empirically
(Appendix~\ref{app:sl})---the signs of the coefficients are valid even
if the magnitudes are attenuated.  The evidence for the directional
claim therefore rests not on any single $p$-value but on the conjunction
of three facts: both coefficients carry their theoretically predicted
signs across all specifications and measurement draws; the shadow
variables improve out-of-sample prediction; and the pattern survives
country fixed effects.

Restricting Stage~2 to the labeled subsample (country-years with
observed intervention coding) yields coefficients close to the
full-sample estimates, confirming that the results are not driven by
out-of-sample extrapolation of the classifier
(\citealt{knoxlucascho2022}, \S5.6).

The fixed-effects specification (Column~3) absorbs all time-invariant
confounders---terrain, ethnic composition, colonial history---and
identifies the shadow effect from within-country temporal variation
alone.  Both coefficients retain their signs: $\gamma^G = -1.74$
(SE $= 0.47$), $\gamma^O = +0.87$ (SE $= 0.47$).  Only 71 of 171
countries exhibit variation in onset and therefore contribute to
identification.  (These standard errors are from maximum likelihood and
do not incorporate the $T \times P$ correction; they should be
interpreted as lower bounds on uncertainty.)

\begin{table}[t]
  \centering
  \begin{tabular}{@{}lrrr@{}}
    \toprule
    & Baseline & Entrants & FE Entrants \\
    \midrule
    Polity2\textsubscript{lag}         & $-$0.026 (0.014) & $-$0.032\textsuperscript{*} (0.014) & $-$0.043\textsuperscript{*} (0.019) \\
    GDP/cap (log)\textsubscript{lag}   & $-$0.670\textsuperscript{***} (0.112) & $-$0.500\textsuperscript{***} (0.122) & $-$0.232 (0.270) \\
    Pop (log)\textsubscript{lag}       & +0.110 (0.061) & +0.095 (0.066) & +0.411 (0.697) \\
    Mountainous                        & +0.223\textsuperscript{***} (0.063) & +0.126 (0.067) & \\
    Noncontiguous                      & +0.532\textsuperscript{*} (0.225) & +0.538\textsuperscript{*} (0.225) & \\
    Oil                                & +0.450\textsuperscript{*} (0.203) & +0.328 (0.208) & \\
    New state                          & +1.003\textsuperscript{**} (0.342) & +1.168\textsuperscript{**} (0.348) & +1.218\textsuperscript{**} (0.385) \\
    Instability\textsubscript{lag}     & +0.567\textsuperscript{**} (0.180) & +0.513\textsuperscript{**} (0.179) & +0.617\textsuperscript{**} (0.196) \\
    Prior war                          & +0.529\textsuperscript{**} (0.195) & +0.450\textsuperscript{*} (0.200) & $-$0.313 (0.203) \\
    Ethnic frac.                       & +0.409 (0.293) & +0.562 (0.299) & \\
    Religious frac.                    & +0.673 (0.392) & +1.225\textsuperscript{**} (0.411) & \\
    Year                               & +0.011\textsuperscript{*} (0.005) & +0.006 (0.005) & $-$0.003 (0.019) \\
    \midrule
    $\hat{E}^G$ (asinh)                &        & $-$1.62 (1.37) & $-$1.74\textsuperscript{***} (0.47) \\
    $\hat{E}^O$ (asinh)                &        & +0.64 (1.10) & +0.87 (0.47) \\
    \midrule
    Country FEs                        & No & No & Yes \\
    $N$                                & 8,792 & 8,792 & 3,750 \\
    Onsets                             & 184 & 184 & 184 \\
    PRL                                & 10.1\% & 11.9\% & \\
    AUC                                & 0.771 & 0.787 & \\
    \bottomrule
  \end{tabular}
  \caption{Logistic regression estimates for civil war onset.  Columns~1--2
    are pooled logit; Column~3 is conditional (fixed-effects) logit.
    Point estimates and standard errors for the shadow variables
    ($\hat{E}^G$, $\hat{E}^O$) in Column~2 are from a $T \times P$ bootstrap
    following \citet{knoxlucascho2022}: $T = 25$
    measurement-model draws $\times$ $P = 200$ pairs-cluster bootstrap
    replications, pooling all 5,000 coefficient vectors.  This accounts
    for uncertainty in both the Stage~1 measurement model and the Stage~2
    primary analysis.  All other standard errors are from maximum
    likelihood.  Column~3 standard errors do not incorporate the $T \times P$
    correction.  Time-invariant covariates are absorbed by country fixed
    effects in Column~3.  \textsuperscript{*}~$p < 0.05$, \textsuperscript{**}~$p < 0.01$, \textsuperscript{***}~$p < 0.001$.}
  \label{tab:coefs}
\end{table}

\subsection{Heterogeneous Utilities}

The learned-proxy approach permits a question that structural designs
cannot easily address: which types of interveners carry disproportionate
weight in the onset calculus?  The Entrants model constrains all
interveners to carry the same utility weight regardless of type.  The
theoretical framework permits the utility
parameters to vary with intervener characteristics: setting
$u_O^i(O) = \gamma_0^O + \gamma_P^O z_i$ and
$u_O^i(G) = \gamma_0^G + \gamma_P^G z_i$ for some intervener attribute
$z_i$ and aggregating produces the Entrants variables plus a pair of
type-weighted interaction terms.  I apply this to major-power status
(\textit{Powers}), geographic contiguity (\textit{Neighbors}), shared dominant ethnic
group (\textit{Coethnics}), colonial relationship (\textit{Rulers}), bilateral rivalry
(\textit{Rivals}), and military balance (\textit{DOE}).  Each type-disaggregated
specification nests the Entrants variables.

\begin{table}[t]
  \centering
  \begin{tabular}{@{}lrrrr@{}}
    \toprule
    Model & IS PRL & IS AUC & OOS PRL & OOS AUC \\
    \midrule
    Baseline       & 10.1\% & 0.771 & $-$1.0\% & 0.650 \\
    Entrants       & 11.9\% & 0.787 & +0.4\%   & 0.672 \\
    Powers         & 14.3\% & 0.809 & +1.7\%   & 0.707 \\
    Neighbors      & 14.5\% & 0.814 & +2.4\%   & 0.710 \\
    Coethnics      & 12.8\% & 0.798 & +0.4\%   & 0.680 \\
    Rulers         & 11.9\% & 0.787 & $-$0.8\% & 0.668 \\
    Rivals (bin)   & 14.4\% & 0.814 & +1.7\%   & 0.706 \\
    Rivals (cts)   & 14.7\% & 0.820 & +2.6\%   & 0.712 \\
    DOE            & 11.9\% & 0.788 & +0.2\%   & 0.671 \\
    Full           & 22.3\% & 0.866 & +4.2\%   & 0.778 \\
    \bottomrule
  \end{tabular}
  \caption{In-sample and out-of-sample fit across model specifications.
    Each extended specification nests the aggregate Entrants variables.
    IS = in-sample; OOS = leave-one-onset-group-out cross-validation.
    PRL = proportional reduction in log-loss relative to the class-frequency null.
    AUC = area under the ROC curve.}
  \label{tab:topfit}
\end{table}

Six of the nine shadow-augmented specifications outperform the Baseline
out-of-sample, with Powers, Neighbors, Rivals~(bin), and
Rivals~(cts) achieving the largest gains (OOS PRL
$\approx 1.7$--$2.6\%$, OOS AUC $\approx 0.71$), indicating that major
powers, contiguous states, and hostile dyads carry disproportionate weight
in shaping onset expectations.  The Full model achieves the highest
discriminative ability (OOS AUC = 0.778) and is the only specification
with OOS PRL above 4\%, though its 26-parameter specification risks
overfitting.

The shadow measure also subsumes the two most direct existing proxies for
anticipated intervention.  On the common sample
(6,973 country-years, 1950--2005), the Lake security hierarchy score
from \citet{cunningham2016} adds no explanatory power beyond the shadow
variables: the hierarchy coefficient is effectively zero
when the Entrants variables are included (Appendix~\ref{app:cunningham},
Figure~\ref{fig:hierarchy}).  On the 7,772 country-years where the
\citet{gibilisco2022} structural P5 intervention probabilities are
available, the pattern is the same: the G\&M aggregate score adds nothing
to the Entrants specification (LR $= 0.06$, $p = 0.80$), but the shadow
variables remain highly significant when added to a model that already
includes the G\&M score (LR $= 26.4$, $p < 0.001$).  The learned proxy
absorbs what both the hierarchy and the structural-game estimates
capture, while adding variation from non-P5 interveners, bilateral
relationships, and the spatial environment that neither alternative
covers.
